@@M0@@
where @@M1@@ is a homogeneous ideal of @@M2@@, of degree @@M3@@. It follows from definition (4.9.1) that @@M4@@ is the separated completion of @@M5@@ for the topology defined by the @@M6@@. This topology is finer than the topology defined by weight; we shall show that it is equivalent to it, which will imply the proposition.

Put @@M7@@, and let @@M8@@ be the @@M9@@-algebra generated by generators @@M10@@ subject to the relations:
@@M11@@
where the @@M12@@ denote the elementary symmetric functions. By 4.4, the canonical homomorphism @@M13@@ is injective. We obtain:
@@M14@@
Since @@M15@@ is a polynomial, so is @@M16@@ for every @@M17@@. Consequently, the @@M18@@ are nilpotent, and there is an integer @@M19@@ such that every polynomial of degree @@M20@@ in the @@M21@@, with coefficients in @@M22@@, is zero. Thus every polynomial of @@M23@@, of weight @@M24@@, is zero, i.e. every polynomial of @@M25@@ of weight @@M26@@ belongs to @@M27@@. It follows that the two topologies under consideration are equivalent, which completes the proof.

\sgaStableTarget{sga6.e06.sec.5}\section*{5. Applications to projective bundles; study of \texorpdfstring{@@M28@@}{f*}}

\sgaStableTarget{sga6.e06.sec.5.1}5.1.
In this section, @@M29@@ is a quasi-compact scheme, @@M30@@ an @@M31@@-Module locally free of rank @@M32@@, and @@M33@@ the projective bundle associated to @@M34@@. On @@M35@@, there is the canonical exact sequence:
@@M36@@
